\documentclass[a4paper,12pt,twoside]{memoir}

% Castellano
\usepackage[spanish,es-tabla]{babel}
\selectlanguage{spanish}
\usepackage[utf8]{inputenc}
\usepackage[T1]{fontenc}
\usepackage{lmodern} % scalable font
\usepackage{microtype}
\usepackage{placeins}
\usepackage{listings}

\RequirePackage{booktabs}
\RequirePackage[table]{xcolor}
\RequirePackage{xtab}
\RequirePackage{multirow}

% Links
\PassOptionsToPackage{hyphens}{url}\usepackage[colorlinks]{hyperref}
\hypersetup{
	allcolors = {red}
}

% Ecuaciones
\usepackage{amsmath}

% Rutas de fichero / paquete
\newcommand{\ruta}[1]{{\sffamily #1}}

% Párrafos
\nonzeroparskip

% Huérfanas y viudas
\widowpenalty100000
\clubpenalty100000

% Evitar solapes en el header
\nouppercaseheads

% Comando para insertar una imagen en un lugar concreto.
% Los parámetros son:
% 1 --> Ruta absoluta/relativa de la figura
% 2 --> Texto a pie de figura
% 3 --> Tamaño en tanto por uno relativo al ancho de página
\usepackage{graphicx}
\newcommand{\imagen}[3]{
	\begin{figure}[!h]
		\centering
		\includegraphics[width=#3\textwidth]{#1}
		\caption{#2}\label{fig:#1}
	\end{figure}
	\FloatBarrier
}

\newcommand{\imagenflotante}[2]{
	\begin{figure}%[!h]
		\centering
		\includegraphics[width=0.9\textwidth]{#1}
		\caption{#2}\label{fig:#1}
	\end{figure}
}



% El comando \figura nos permite insertar figuras comodamente, y utilizando
% siempre el mismo formato. Los parametros son:
% 1 -> Porcentaje del ancho de página que ocupará la figura (de 0 a 1)
% 2 --> Fichero de la imagen
% 3 --> Texto a pie de imagen
% 4 --> Etiqueta (label) para referencias
% 5 --> Opciones que queramos pasarle al \includegraphics
% 6 --> Opciones de posicionamiento a pasarle a \begin{figure}
\newcommand{\figuraConPosicion}[6]{%
  \setlength{\anchoFloat}{#1\textwidth}%
  \addtolength{\anchoFloat}{-4\fboxsep}%
  \setlength{\anchoFigura}{\anchoFloat}%
  \begin{figure}[#6]
    \begin{center}%
      \Ovalbox{%
        \begin{minipage}{\anchoFloat}%
          \begin{center}%
            \includegraphics[width=\anchoFigura,#5]{#2}%
            \caption{#3}%
            \label{#4}%
          \end{center}%
        \end{minipage}
      }%
    \end{center}%
  \end{figure}%
}

%
% Comando para incluir imágenes en formato apaisado (sin marco).
\newcommand{\figuraApaisadaSinMarco}[5]{%
  \begin{figure}%
    \begin{center}%
    \includegraphics[angle=90,height=#1\textheight,#5]{#2}%
    \caption{#3}%
    \label{#4}%
    \end{center}%
  \end{figure}%
}
% Para las tablas
\newcommand{\otoprule}{\midrule [\heavyrulewidth]}
%
% Nuevo comando para tablas pequeñas (menos de una página).
\newcommand{\tablaSmall}[5]{%
 \begin{table}
  \begin{center}
   \rowcolors {2}{gray!35}{}
   \begin{tabular}{#2}
    \toprule
    #4
    \otoprule
    #5
    \bottomrule
   \end{tabular}
   \caption{#1}
   \label{tabla:#3}
  \end{center}
 \end{table}
}

%
%Para el float H de tablaSmallSinColores
\usepackage{float}

%
% Nuevo comando para tablas pequeñas (menos de una página).
\newcommand{\tablaSmallSinColores}[5]{%
 \begin{table}[H]
  \begin{center}
   \begin{tabular}{#2}
    \toprule
    #4
    \otoprule
    #5
    \bottomrule
   \end{tabular}
   \caption{#1}
   \label{tabla:#3}
  \end{center}
 \end{table}
}

\newcommand{\tablaApaisadaSmall}[5]{%
\begin{landscape}
  \begin{table}
   \begin{center}
    \rowcolors {2}{gray!35}{}
    \begin{tabular}{#2}
     \toprule
     #4
     \otoprule
     #5
     \bottomrule
    \end{tabular}
    \caption{#1}
    \label{tabla:#3}
   \end{center}
  \end{table}
\end{landscape}
}

%
% Nuevo comando para tablas grandes con cabecera y filas alternas coloreadas en gris.
\newcommand{\tabla}[6]{%
  \begin{center}
    \tablefirsthead{
      \toprule
      #5
      \otoprule
    }
    \tablehead{
      \multicolumn{#3}{l}{\small\sl continúa desde la página anterior}\\
      \toprule
      #5
      \otoprule
    }
    \tabletail{
      \hline
      \multicolumn{#3}{r}{\small\sl continúa en la página siguiente}\\
    }
    \tablelasttail{
      \hline
    }
    \bottomcaption{#1}
    \rowcolors {2}{gray!35}{}
    \begin{xtabular}{#2}
      #6
      \bottomrule
    \end{xtabular}
    \label{tabla:#4}
  \end{center}
}

%
% Nuevo comando para tablas grandes con cabecera.
\newcommand{\tablaSinColores}[6]{%
  \begin{center}
    \tablefirsthead{
      \toprule
      #5
      \otoprule
    }
    \tablehead{
      \multicolumn{#3}{l}{\small\sl continúa desde la página anterior}\\
      \toprule
      #5
      \otoprule
    }
    \tabletail{
      \hline
      \multicolumn{#3}{r}{\small\sl continúa en la página siguiente}\\
    }
    \tablelasttail{
      \hline
    }
    \bottomcaption{#1}
    \begin{xtabular}{#2}
      #6
      \bottomrule
    \end{xtabular}
    \label{tabla:#4}
  \end{center}
}

%
% Nuevo comando para tablas grandes sin cabecera.
\newcommand{\tablaSinCabecera}[5]{%
  \begin{center}
    \tablefirsthead{
      \toprule
    }
    \tablehead{
      \multicolumn{#3}{l}{\small\sl continúa desde la página anterior}\\
      \hline
    }
    \tabletail{
      \hline
      \multicolumn{#3}{r}{\small\sl continúa en la página siguiente}\\
    }
    \tablelasttail{
      \hline
    }
    \bottomcaption{#1}
  \begin{xtabular}{#2}
    #5
   \bottomrule
  \end{xtabular}
  \label{tabla:#4}
  \end{center}
}



\definecolor{cgoLight}{HTML}{EEEEEE}
\definecolor{cgoExtralight}{HTML}{FFFFFF}

%
% Nuevo comando para tablas grandes sin cabecera.
\newcommand{\tablaSinCabeceraConBandas}[5]{%
  \begin{center}
    \tablefirsthead{
      \toprule
    }
    \tablehead{
      \multicolumn{#3}{l}{\small\sl continúa desde la página anterior}\\
      \hline
    }
    \tabletail{
      \hline
      \multicolumn{#3}{r}{\small\sl continúa en la página siguiente}\\
    }
    \tablelasttail{
      \hline
    }
    \bottomcaption{#1}
    \rowcolors[]{1}{cgoExtralight}{cgoLight}

  \begin{xtabular}{#2}
    #5
   \bottomrule
  \end{xtabular}
  \label{tabla:#4}
  \end{center}
}




\graphicspath{ {./img/} }

% Capítulos
\chapterstyle{bianchi}
\newcommand{\capitulo}[2]{
	\setcounter{chapter}{#1}
	\setcounter{section}{0}
	\setcounter{figure}{0}
	\setcounter{table}{0}
	\chapter*{#2}
	\addcontentsline{toc}{chapter}{#2}
	\markboth{#2}{#2}
}

% Apéndices
\renewcommand{\appendixname}{Apéndice}
\renewcommand*\cftappendixname{\appendixname}

\newcommand{\apendice}[1]{
	%\renewcommand{\thechapter}{A}
	\chapter{#1}
}

\renewcommand*\cftappendixname{\appendixname\ }

% Formato de portada
\makeatletter
\usepackage{xcolor}
\newcommand{\tutor}[1]{\def\@tutor{#1}}
\newcommand{\course}[1]{\def\@course{#1}}
\definecolor{cpardoBox}{HTML}{E6E6FF}
\def\maketitle{
  \null
  \thispagestyle{empty}
  % Cabecera ----------------
\noindent\includegraphics[width=\textwidth]{cabecera}\vspace{1cm}%
  \vfill
  % Título proyecto y escudo informática ----------------
  \colorbox{cpardoBox}{%
    \begin{minipage}{.8\textwidth}
      \vspace{.5cm}\Large
      \begin{center}
      \textbf{TFG del Grado en Ingeniería Informática}\vspace{.6cm}\\
      \textbf{\LARGE\@title{}}
      \end{center}
      \vspace{.2cm}
    \end{minipage}

  }%
  \hfill\begin{minipage}{.20\textwidth}
    \includegraphics[width=\textwidth]{escudoInfor}
  \end{minipage}
  \vfill
  % Datos de alumno, curso y tutores ------------------
  \begin{center}%
  {%
    \noindent\LARGE
    Presentado por \@author{}\\ 
    en Universidad de Burgos --- \@date{}\\
    Tutores: \@tutor{}\\
  }%
  \end{center}%
  \null
  \cleardoublepage
  }
\makeatother


% Datos de portada
\title{Análisis de datos de temperatura y humedad de suelo procedentes de sensores IoT desplegados en un viñedo  \\ Documentación Técnica}
\author{Gabriel Hernández Vallejo}
\tutor{Rubén Ruiz González, Alejandro Merino Gómez}
\date{\today}

\begin{document}

\maketitle



\cleardoublepage



%%%%%%%%%%%%%%%%%%%%%%%%%%%%%%%%%%%%%%%%%%%%%%%%%%%%%%%%%%%%%%%%%%%%%%%%%%%%%%%%%%%%%%%%



\frontmatter


\clearpage

% Indices
\tableofcontents

\clearpage

\listoffigures

\clearpage

\listoftables

\clearpage

\mainmatter

\appendix

\apendice{Introducción}
Al no tratarse de un proyecto de desarrollo de software, sino que se encuentra enmarcado 
dentro del campo de \textit{Machine Learning}, se seguirá una estructura en los anexos algo diferente a la propuesta original.
Se mantienen los apartados clave como son el plan de proyecto, el manual de programador y usuario y se sustituyen el resto
por las diferentes fases de descubrimiento de conocimiento en bases de datos (\textit{KDD}).

Estas se corresponderán con el conocimiento del negocio, conocimiento de datos, preparación de datos, modelado y evalauación
de los resultados.

\apendice{Plan de Proyecto}

\section{Introducción}
En la sección se introduce la fase de planificación del proyecto, tanto la planificación temporal, como
la viabilidad económica y legal.

\section{Planificación temporal}
La planificación del desarrollo del proyecto se realizó siguiendo la metodología ágil \textit{Scrum}, aunque no en su definición completa, puesto
que al ser una única persona en el desarrollo, no permite realizar todos los procedimientos recogidos dentro de los manuales de gestión de proyectos
con metodologías ágiles \cite{book:palacio2021}, como las reuniones diarias ``daily'' para recoger diferentes cuestiones como cúales son los 
problemas que ralentizan el proceso de desarrollo del proyecto.

Por otro lado, la herramienta software que apoyaba su aplicación (ZenHub \cite{misc:zenhub2023}) dejó de estar disponible de forma abierta, de forma
que no pudo seguirse con el modelo \textit{Canvas} que aplicaba la metodología recogiendo en un tablero con varias columnas, según fuera su estado, las 
tareas por cada uno de los \textit{sprints}.

Aun así, se ha continuado aplicando la filosofía ágil en líneas generales:
\begin{itemize}
    \item Se continuó con el desarrollo iterativo incremental mediante los \textit{sprints}.
    \item La duración media de estas iteraciones fueron de dos semanas al comienzo del proyecto y durante
        la mayor parte de su realización, reduciendo el tiempo al final de este.
    \item En la terminación de los \textit{sprints} se realizaba un incremento, manteniendo reuniones con los 
        tutores para la planificación de la iteración entrante y la revisión de posibles errores en el desarrollo.
    \item Las tareas surgidas de la reunión se creaba, estimaban y priorizaban, en un inicio añadiéndolas al tablero
        \textit{Canvas}.
\end{itemize}

La estimación fue realizada empleando los \textit{story points} disponibles en ZenHub, de forma que estos iban de las 
tareas más sencillas de implementar y rápidas, a las más complejas y que requerían de mayor tiempo
de desarrollo.

\subsection{Sprint 0 (18/01/2023 - 31/01/2023)}
En este \textit{sprint} se presentó el proyecto, indicándose los primeros objetivos, lo que se enmarca en la fase de comprensión
del negocio.
Las primeras tareas se correspondían con la fase de comprensión de datos y a su vez la preparación de los mismos, 
puesto que el conjunto proporcionado contaba con diferentes casuísticas de las que en un comienzo se tenía constancia.

Entre los problemas encontrados previos al análisis en profundidad de los datos se encontraban lecturas incorrectas (valores erróneos) 
debido a problemas de \textit{overflow} de variables en los controladores de cada sensor, de forma que cuando la batería bajaba
de los 3.1V, las lecturas en el conjunto de datos se trataban de errores en su mayoría.

Además, en ciertos periodos de tiempo los sensores habían estado desconectados, por lo que se produjeron lecturas nulas (valores faltantes)
que necesitaban ser tratadas.

Por otro lado, las lecturas en el pluviómetro instalado en las dependencias del viñedo en algunas de las muestras arrojaban valores
inverosímiles debido a la forma en la que este se encontraba instalado, de manera que el balancín empleado para medir las diferencias de 
precipitaciones entre instantes se activaba artificialmente.

Durante el desarrollo de la iteración se comenzó con el tratamiento de los datos nulos, tomando la decisión de eliminar las entradas, pues se
contaba con un número elevado de muestras y realizar la recuperación de estas no parecía viable debido a la naturaleza de los datos.
Por otro lado, para los problemas relacionados con el pluviómetro se solicitaron claves \textit{API} para poder acceder a los registros 
de estaciones meteorológicas cercanas a la zona y poder comprobar los registros.

Por otro lado, se comenzó a realizar las gráficas de los datos para intentar encontrar correlaciones en los datos y poder hallar diferentes 
errores en el conjunto de datos.

\subsection{Sprint 1 (31/01/2023 - 14/02/2023)}
En este \textit{sprint} se presentaron los avances de la primera iteración y se acordó la continuación del proceso de comprensión de datos y 
preparación de los mismos.

Los esfuerzos se centraron entonces en modificar las gráficas de visualización de datos para continuar con la inspección visual de los sensores 
y comenzar con los datos del pluviómetro.
Con estas inspecciones se consigue encontrar anomalías en ciertos sensores como en el caso del segundo y cuarto.

\subsection{Sprint 2 (14/02/2023 - 28/02/2023)}
En la reunión de la iteración se propusieron diferentes objetivos:
\begin{itemize}
    \item En el sensor 2 eliminar con variación excesiva en temperatura, corrigiendo los datos ruidosos.
    \item En el sensor 3 se encontró que el salto de humedad no era genuino.
    \item En el sensor 4 intentar realizar una media móvil para arreglar variables como la temperatura al emplear variaciones diarias en lugar de 
        muestreos cada 5 minutos.
\end{itemize}

Durante el desarrollo del \textit{sprint} se continuó con la selección de datos, encontrando que el efecto de la lluvia provocaba ciertos cambios 
bruscos en los datos como en el caso del sensor 5, pero estando algunos aparentemente no relacionados con este fenómeno.
Se comenzó, por otro lado con la selección y limpieza de los datos mediante una columna adicional de validez en los ficheros correspondientes.

Además, se implementó la detección de \textit{Outliers} mediante el rango intercuartílico, para eliminar el posible ruido del conjunto de datos
de los sensores.

\subsection{Sprint 3 (28/02/2023 - 14/03/2023)}
En la reunión de la iteración se acordó realizar la recuperación de los datos del sensor 8 hasta la variabilidad excesiva de los valores de las
variables observadas, así como el intento de mejorar la calidad del conjunto de datos en otros sensores realizando procesos similares.
Por otro lado, se propuso emplear WeatherBit como \textit{API} meteorológica para eliminar los datos que más distaran de las lecturas del 
pluviómetro desplegado.

Durante el desarrollo del \textit{sprint} se modificaron los umbrales empleados en la detección y manejo de valores extremos, para eliminar la mayoría
cantidad de ruido posible en todos los sensores.
Por otro lado, se fueron finalizando las selecciones de datos de los sensores, de forma que se realizó la recuperación de los datos en sensores cuyas 
variabilidades se debían a factores externos (fueron sacados de sus posiciones originales).

En lo referente del sensor 8 se llegó a la conclusión de que buena parte de los datos eran irrecuperables debido a las variabilidades aleatorias 
con las que parecía contar.

Por otro lado, se comenzó con la limpieza de los datos de los sensores empleando los datos de varias \textit{API}, en primer lugar se trató de emplear
la mencionada anteriormente, siendo el \textit{endpoint} ideal de licencia de pago.
Probando con OpenWeatherMap se presentó el caso al soporte, ampliando la licencia de usuario para poder acceder a l utilidad requerida, sin embargo no
se permitía retrotraerse en más de un año en las lecturas registradas.

Finalmente se decide emplear los servicios de la AEMET, debido principalmente a que permitía obtener los datos de hacía más de un año, sin embargo, contaba
con inconvenientes, y es que no permitía obtener datos por hora, sino por día y la estación meteorológica más cercana se encontraba a unos 15 Kilómetros
del despliegue.

\subsection{Sprint 4 (14/03/2023 - 28/03/2023)}
En la reunión del \textit{sprint} se acuerda comenzar con la fase del modelado a la vista de la aparente correcta selección y limpieza de los datos 
proporcionados.
Se plantea la creación de una matriz de correlación para observar cúales de los atributos seleccionar para la creación del modelo.

Durante el desarrollo de la iteración los esfuerzos se centran en la creación del entorno virtual con las dependencias correspondientes que permitan 
realizar los modelados, además de dejar el fichero de pre-procesamiento de datos comentado con las explciaciones pertinentes.

Por otro lado, se realizan avances sustanciales en la memoria del proyecto añadiendo información el la introducción, objetivos y técnicas y herramientas,
además de realizar una reestructuración de ficheros y una modificación de la forma en la que se realizan las gráficas de los datos.

Se comienza con la creación tentativa del modelo neuronal en busca de conocimiento del dominio, que permita profundizar en los módulos empleados para tal
propósito, realizando regresión simple.
Se investiga sobre las redes neruonales recurrentes, más concretamente sobre la aplicación de modelos como \textit{GRU} y \textit{LSTM} en problemas
similares (predicciones de tiempo atmosférico).

\subsection{Sprint 5 (28/03/2023 - 11/04/2023)}
En la reunión de la iteración se muestran los avances en memoria y en la creación de modelos, por otro lado, se plantea la forma de realizar la regresión
empleando Keras.

Durante el \textit{sprint} se modifica la manera en la que se crean los datos de entrada a los modelos, empleando medias diarias para tal propósito y 
dividiendo los datos en dos conjuntos: entrenamiento y validación. Además, se crean las gráficas de dispersión que permiten comparar visualmente
las predicciones y valores reales.

Por otro lado, se continúa con el desarrollo de la memoria, añadiendo conceptos teóricos sobre el proceso de extracción de conocimiento de bases de datos 
(\textit{KDD}), así como modificando algunos apartados.

\subsection{Sprint 6 (11/04/2023 - 25/04/2023)}
En la planificación se considera entrenar a los modelos neuronales únicamente con valores adecuados, pues hasta ahora pueden existir saltos temporales que
pueden afectar a la regresión, realizando medias temporales cada cierto periodo de tiempo para eliminar los posibles problemas de \textit{offset} de los 
muestreos originales.

En el desarrollo del \textit{sprint}, se modificó la selección en los datos de entrada a los modelos, de forma que se realizaban las medias diarias de los
diferentes atributos y se les aplicaba el filtro de Hodrick-Presscott, un filtro de tendencias, que permite obtener las tendencias en una serie temporal,
de manera que los modelos sean capaces de generalizar de forma más adecuada.

Por otro lado, se continuó con la memoria, en concreto con los conceptos teóricos relacionados con la fase de modelado.

\subsection{Sprint 7 (25/04/2023 - 09/05/2023)}
En la reunión de la iteración se revisan los cambios realizados, llegando a la consideración de emplear varios dias previos como entrada a los modelos,
utilizando una ventana deslizante para seleccionar datos, de manera que se eviten los saltos temporales y estudiar la incorporación de una componente temporal 
como entrada en la regresión.

De esta forma, en el desarrollo se aplicó la ventana deslizante para la selección de datos de entrada en el modelo, de manera que se pudieran evitar los saltos
temporales.

\subsection{Sprint 8 (09/05/2023 - 16/05/2023)}
En la planificación se acordó intentar rebajar los periodos de tiempo, y que se emplearan datos medios por hora en lugar de diarios y probar los hiperparámetros
de las redes, intentando variar las neuronas por capa, el número de estas, el ratio de aprendizaje, etc. para observar el comportamiento de los modelos y
determinar tanto si continúan generalizando como si es necesario realizar cambios en el proceso de conocimiento de datos y preparación de estos.

De esta manera, se redujo la frecuencia de las medias de los datos, realizando la integración de los datos de los sensores en un único fichero, para poder
entrenar los modelos empleando un único conjunto. Esto obligó a cambiar la forma en la que se graficaban los resultados, mostrando las gráficas de 
dispersión por sensor y atributo.

Por otro lado, se dividió el conjunto de datos en tres subconjuntos diferentes: un conjunto para entrenamiento, otro para validación y el restante para test, 
permitiendo la parametrización en la división de estos.

\subsection{Sprint 9 (16/05/2023 - 23/05/2023)}
En la reunión de la iteración se acordó tratar de predecir tiempos más largos en lugar de un único instante como realizaban los modelos hasta el momento;
previsiblemente el error sería mayor, por lo que se mencionó permitir parametrizar este valor.

Ante el objetivo principal planteado hubo que realizar diferentes modificaciones en la selección de los datos de salida y los parámetros, penalizando el rendimiento
general de los modelos.

Se avanzó en la memoria en los conceptos teóricos y se realizaron algunas correcciones.

\subsection{Sprint 10 (23/05/2023 - 30/05/2023)}
En la planificación básicamente se indicó la continuación en los esfuerzos de realización de memoria y anexos.
De esta manera, este \textit{sprint} se centró en la finalización de los mismos.

\section{Estudio de viabilidad}
En esta sección se desglosarán la viabilidad económica y legal del proyecto, en cuanto a la primera indicará los costes derivados del desarrollo en un enotorno real.
En lo referente la segunda, se presentarán las licencias empleadas en el proyecto y su implicación con librerías de terceros.

\subsection{Viabilidad económica}
En términos de viabilidad económoca es necesario hacer una diferenciación entre los costes y los beneficios que lleva realizar el proyecto.

\subsubsection{Costes}
Los costes que pueden surgir del proyecto en un entorno empresarial pueden desglosarse en los siguientes:

\textbf{Costes de personal:}

El desarrollo se ha realizado con un única persona empleada a tiempo completo en aproximadamente 4 meses, de forma que se consideran los siguientes 
costes:
\tablaSmallSinColores{Costes de Hardware}{l l}{Coste Hardware}{\textbf{Concepto} & \textbf{Coste} \\}
{Salario mensual neto & 1.080€ \\
Retención I.R.P.F(12\%) & 185,46€ \\
Seguridad Social(31,1\%) & 480,64€ \\
\hline
Salario mensual bruto & 1.545€ \\
\hline
\textbf{Total 4 meses} & 6.181,92€ \\
}

El porcentaje mensual de aporte a la Segurida Social se calcula como el 0,2\% al Fondo de Garantía Salarial (FOGASA), más el 0,6\%, el 6,7\% de
la prestación por desempleo y el 23,6\% de contingencias comunes \cite{misc:ss2023}.

\textbf{Costes de hardware:}

El hardware empleado tiene un tiempo de amortización aproximado de 4 años.
\tablaSmallSinColores{Costes de Hardware}{l l l}{Coste Hardware}{\textbf{Concepto} & \textbf{Coste} & \textbf{Coste amortizado} \\}
{Ordenador portátil & 900€ & 45€ \\
\hline
\textbf{Total} & 900€ & 45€ \\
}

\textbf{Costes de software:}

En cuanto a los costes de software, hay ciertos programas o herramientas empleadas durante el desarrollo que requieren de licencia de pago,
estos contarán con un tiempo de amortización estimado de 2 años.
\tablaSmallSinColores{Costes de Software}{l l l}{Coste Software}{\textbf{Concepto} & \textbf{Coste} & \textbf{Coste amortizado} \\}
{Windows 10 Home & 148€ & 18,5€ \\
\hline
\textbf{Total} & 148€ & 18,5€ \\
}

\textbf{Costes varios:}

Al igual que en otras situaciones empresariales, en los desarrollos de software surgen tanto costes inesperados, como fijos.
Se reflejarán aquellos costes variados que aparentemente se tratarían de costes fijos.
\tablaSmallSinColores{Costes varios}{l l}{Coste varios}{\textbf{Concepto} & \textbf{Coste} \\}
{Memoria y anexos & 50€ \\
Alquiler de espacio de trabajo & 592€ \\
Internet & 120€ \\
\hline
\textbf{Total} & 762€ \\
}

\textbf{Costes totales:}

Los costes del desarrollo del proyecto son los siguientes:
\tablaSmallSinColores{Costes totales}{l l}{Coste totales}{\textbf{Concepto} & \textbf{Coste} \\}
{Personal & 6.181,92€ \\
Hardware & 900€ \\
Software & 148€ \\
Varios & 762€ \\
\hline
\textbf{Total} & 7.991,92€ \\
}

\subsubsection{Beneficios}
En cuanto a los beneficios generados, en el \textit{scope} del proyecto no había cabida a un despliegue en los meses de desarrollo de los 
modelos, por lo que a corto plazo no habría ningún beneficio aparente.

Sin embargo, la forma de obtener los beneficios entraría dentro de una segunda fase, cuando se realice el despliegue del resultados
en algún producto.

\subsection{Viabilidad legal}
En esta sección se desglosarán los temas relacionados con las licencias de uso de los productos software.

En primer lugar se analizará la licencia más conveniente en el desarrollo del proyecto, teniendo en cuenta las dependencias empleadas y 
sus respectivas licencias, puesto que la selección estará limitada por las restricciones de estas principalmente.

\tablaSmallSinColores{Licencias de las dependencias utilizadas}{m{6em} m{5em} m{12em} m{5em}}{Licencias dependencias}{\textbf{Dependencia} & \textbf{Versión} & \textbf{Descripción} & \textbf{Licencia}\\}
{TensorFlow & 2.11.0 & Biblioteca de aprendizaje automático & Apache v2.0 \\
Pandas & 1.5.3 & Biblioteca especializada en manipulación y análisis de datos & BSD \\
Matplotlib & 2.11.0 & Biblioteca para generación de gráficos & BSD \\
IPython Kernel for Jupyter & 6.21.3 & Biblioteca para manipulación de Python Notebooks & BSD \\
Statsmodel & 0.13.5 & Biblioteca con modelos y funciones estadísticas & BSD \\
}

De esta forma, se debe seleccionar una licencia compatible con Apache v2.0 y BSD, siendo la primera la más restrictiva de las licencias, que requiere la conservación del aviso de derecho
de autor y un descargo de responsabilidad, sin embargo, no es \textit{copyleft} \cite{misc:apache2}.

Por su versatilidad se escoge una licencia MIT, de forma que el se trata de una licencia permisiva de software libre, imponiendo muy poocas restricciones y otorgando una
buena compatibilidad con otras licencias.

\apendice{Comprensión del negocio}

\section{Introducción}
En la fase de comprensión del negocio se pretende analizar los objetivos y requisitos del 
proyecto.

Estos quedarán bien marcados por las necesidades establecidas en las descripciones
del TFG, quedando reflejadas en la sección de introducción de la memoria.

A continuación se resumen tanto los requisitos como los objetivos.

\section{Objetivos del proyecto}
Entre los principales objetivos del proyecto se encuentran la realización de un análisis
de los datos proporcionados de forma que se pueda facilitar la comprensión de los
datos recogidos mediante representaciones gráficas, buscando, de esta forma correlaciones
y para poder encontrar modelos predictivos para anticipar la evolución futura de las 
diferentes variables en un plazo de tiempo determinado.

\section{Requisitos del proyecto}
Entre los rquisitos del proyecto se encontraría la creación de modelos adecuados a los datos
que predigan correctamente las diferentes variables objetivo.

\apendice{Comprensión de los datos}

\section{Introducción}
En la etapa de comprensión de los datos se creará el conjunto inicial y se comprobará
si este es adecuado, para, en caso contrario, poder continuar recopilando.

Los plazos del proyecto impiden realizar esta fase en toda su extensión, puesto que
a pesar de ser el proceso de análisis de datos un proceso iterativo, se encuentra 
sujeto a unos plazos inamovibles.

\section{Descripción de los datos}
Se proporcionan un conjunto de datos compuesto por varios sensores (en concreto 8 
sensores \textit{IoT}) y un pluviómetro desplegados en un viñedo.

Los datos de los sensores proporcionados cuentan con los siguientes atributos:

\tablaSmallSinColores{Atributos de los datos: sensores}{m{5em} m{15em} m{10em}}{Atributos sensores}{\textbf{Atributo} & \textbf{Descripción}  & \textbf{Medida} \\}
{
ts          & \textit{Timestamp} de la muestra                    & ms \\
fecha       & Fecha de la muestra                                 & Fecha en formato yy/mm/dd, hh:mm:ss \\
batería     & Nivel de batería del sensor                         & V \\
t\_ext      & Temperatura ambiental                               & ºC \\
h\_ext      & Humedad ambiental                                   & \% \\
t\_C\_cal   & Temperatura superficial a 20cm.                     & ºC \\
h\_C\_cal   & Humedad superficial a 20cm.                         & \% \\
t\_L\_cal   & Temperatura a profundidad mayor a t\_C\_cal         & ºC \\
h\_L\_cal   & Humedad superficial a profundiad mayor a h\_C\_cal  & \% \\
h\_C        & Humedad superficial sin calibrar                    & Valor relacionado con capacitancias \\
h\_L        & Humedad profunda sin calibrar                       & Valor relacionado con capacitancias \\
}

Los datos del pluviómetro cuentan con los siguientes atributos:
\tablaSmallSinColores{Atributos de los datos: pluviómetro}{m{7em} m{13em} m{10em}}{Atributos pluviómetro}{\textbf{Atributo} & \textbf{Descripción}  & \textbf{Medida} \\}
{
ts              & \textit{Timestamp} de la muestra                      & ms \\
fecha           & Fecha de la muestra                                   & Fecha en formato yy/mm/dd, hh:mm:ss \\
batería         & Nivel de batería del sensor                           & V \\
pluv            & Contador del incremento de activación del balancín    & N/A \\
pluv\_delta     & Incremento de activación del balancín                 & N/A \\
pluv\_deltaMM   & Precipitaciones registradas desde la última lectura   & litros/\(m^2\) \\
}

\section{Visualización de los datos}
Para dar una idea del estado de los datos se proporcionan unas gráficas de dispersión
con las lecturas realizadas en cada uno de los sensores y el pluviómetro.
En estas podrán observarse los valores \textbf{sin procesar} de cada atributo 
en función del tiempo.

En el caso de los sensores, los atributos \textit{h\_C} y \textit{h\_L} se obviarán debido
a que se tratan de datos redundantes con las respectivas humedades calibradas. 
Además, los datos de batería y fecha tampoco se emplearán en la creación de las gráficas,
puesto que no aportan información útil para los objetivos del proyecto más allá de posibles
errores de muestreo.

\subsection{Sensor 1}

% Sensor 1
\imagen{raw_data/sensor1.png}{Datos no procesados: sensor 1}{1}

En el primer sensor, como puede observarse en la Figura \ref{fig:raw_data/sensor1.png},
se encuentran varias situaciones que a \textit{priori} son sencillas de solucionar, 
y es que existe una cantidad importante de muestras para las que no contamos con datos,
además de existir valores claramente inverosímiles (como en el caso de la humedad
más profunda, o el aumento tan repentino de la humedad superficial).
A pesar de esto, el resto parecen estar dentro de la normalidad, por lo que
el trabajo en el sensor mencionado será uno de los menos tediosos.

Las subidas tan repentinas como las observadas en la imágen posteriormente se demostró
que eran debidas al agotamiento de la batería del respetivo sensor, que poco antes
de producirse arrojaban lecturas artificiales para luego realizar lecturas nulas.

\subsection{Sensor 2}

% Sensor 2
\imagen{raw_data/sensor2.png}{Datos no procesados: sensor 2}{1}

En el segundo caso, además de las lecturas para las que no contemos con valores,
el la Figura \ref{fig:raw_data/sensor2.png} se observa cómo llegado un instante
las sondas de temperatura y humedad superficiales comienzan a realizar lecturas
muy variadas a lo largo del tiempo, para posteriormente dejar de producir datos.

Además, puede observarse la presencia de cierto ruido durante los muestreos previos
al suceso mencionado en las sondas de humedad, lo que posiblemente sea debido a 
las capacitancias de estas.

% Una investigación posterior confirmó que el sensor había sido desplazado de su 
% posición original debido probablemente a las labores realizadas en los viñedos.

\newpage
\subsection{Sensor 3}

% Sensor 3
\imagen{raw_data/sensor3.png}{Datos no procesados: sensor 3}{1}

En el este sensor, como puede observarse en la Figura \ref{fig:raw_data/sensor3.png},
además de los muestras para los que no tenemos valores (que será la tónica habitual
en el conjunto de datos), existen unas subidas remarcables de humedades, con unas
bajadas que pueden considerarse aceleradas o repentinas.

La mayor parte de estas son genuinas y debidas a las lluvias contrastadas esos días, sin 
embargo, la última de las subidas se observa contextualmente errónea, provocando 
un cambio brusco en ambas humedades del suelo.

\newpage
\subsection{Sensor 4}

% Sensor 4
\imagen{raw_data/sensor4.png}{Datos no procesados: sensor 4}{1}

En este caso, pueden observarse en la Figura \ref{fig:raw_data/sensor4.png}, diferentes
situaciones. La primera de ellas, como en el resto de sensores son los valores faltantes; la
segunda es el ruido existente en ciertos periodos de tiempo.

Tras este ruido instantáneo en los datos, puede observarse cómo la variabilidad de las muestras
aumenta de forma homogénea en todas y cada una de las sondas, a salvedad de las que recogen
datos ambientales.

Las posteriores investigaciones indicaron que de nuevo el sensor había sido desplazado
de su situación inicial.

\newpage
\subsection{Sensor 5}

% Sensor 5
\imagen{raw_data/sensor5.png}{Datos no procesados: sensor 5}{1}

En la Figura \ref{fig:raw_data/sensor5.png} puede observarse diferentes situaciones.
Una de ellas es la repetida falta de valores para algunos instántes de tiempo; la 
segunda es las aparentes subidas repentinas de humedades.

Estas siguen el mismo patrón en todas las situaciones, aumento que puede clasificarse
como repentino, y un descenso moderado a lo largo del tiempo.
De esta forma, todas se contrastaron como genuinas observando ciertas situaciones como
el tiempo atmosférico en los instántes de subida.

Este sensor se trata de uno con los datos más limpios de forma predefinida de todo
el conjunto de datos, sino el que en mejor estado se encuentra.

\newpage
\subsection{Sensor 6}

% Sensor 6
\imagen{raw_data/sensor6.png}{Datos no procesados: sensor 6}{1}

En el conjunto de datos que se puede observar en la Figura \ref{fig:raw_data/sensor6.png},
al igual que en los anteriores sensores, pueden observarse valores faltantes, además de 
otras situaciones.

Una de ella es la presencia generalizada de ruido en el conjunto de ejemplos en las lecturas
de las sondas que miden variables del suelo.

Por otro lado, en un instante de tiempo las humedades aumentan de forma repentina, posiblemente
por la lluvia, sin embargo en la humedad del suelo a mayor profundidad, hay una variabilidad
que se consideró como irrecuperable por su estado.

\newpage
\subsection{Sensor 7}

% Sensor 7
\imagen{raw_data/sensor7.png}{Datos no procesados: sensor 7}{1}

En este caso pueden observarse en la Figura \ref{fig:raw_data/sensor7.png} diferentes
situaciones, además de la falta ya mencionada en el resto de sensores de valores
en algunas muestras.

En primer lugar, en las sondas de humedad no ambientales se comienza a producir un aumento
que podría considerarse como exponencial en el tiempo debido al agotamiento de la batería.

En segundo lugar, como en el sensor 4, se produce un aumento significativo de la variabilidad
de la temperatura del suelo, produciéndose, además, un aumento de la sonda de
humedad más profunda.

De igual forma que en el sensor mencionado, las labores en los cultivos habían producido un
cambio en el emplazamiento original del sensor, lo que producía estas situaciones mencionadas.
Con respecto a la falta de los datos de humedad y temperatura más superficiales se 
comprobará que fue debida a la rotura de las respectivas sondas.

\newpage
\subsection{Sensor 8}

% Sensor 8
\imagen{raw_data/sensor8.png}{Datos no procesados: sensor 8}{1}

Como puede observarse en la figura \ref{fig:raw_data/sensor8.png}, en este caso existen 
múltiples problemas en el conjunto de datos.

En primer lugar, los ya mencionados valores faltantes en diferentes instantes de tiempo.
En segundo lugar, los aumentos de temperatura debidos al agotamiento de la batería y en tercer
lugar la alta variabilidad en las temperaturas, que por otro lado, se asemejan en gran medida,
lo que indica que de nuevo el sensor había sido desplazado de su emplazamiento original.

Este será uno de los sensores con el conjunto de datos con más problemas de todos los 
expuestos anteriormente.

\newpage
\subsection{Pluviómetro}

% Pluviómetro
\imagen{raw_data/pluviometro.png}{Datos no procesados: pluviómetro}{1}

En el caso del pluviómetro, los datos estarán sujetos a muestras erróneas debido a la forma
de fijación del mismo, puesto que con las ráfagas de viento se producían oscilaciones
en el poste al que se encontraba anclado, de forma que se realizaban lecturas incorrectas
de las precipitaciones.

\apendice{Preparación de los datos}

\section{Introducción}
En la fase de preparación de los datos se realiza la limpieza de estos para 
que posteriormente puedan ser empleados por los modelos y obtener resultados adecuados
para ser utilizados en un posterior despliegue.

Este proceso se compone de la detección y el tratamiento de varias casuísticas en 
los datos, como los valores faltantes, que se tratan de valores de los que no tenemos
registros por diferentes motivos y los valores extremos, que son aquellos
estadísticamente lejanos al resto, lo que no quiere decir que sean erróneos.

Por otro lado, también es posible encontrar valores erróneos, que pueden ser 
estadísticamente correctos, pero no tener sentido contextual. En este
caso concreto se tratarán de incrementos o bajadas inverosímiles de temperaturas 
y/o humedades del suelo, que se encuentran dentro de los valores naturales,
pero que es difícil que tengan lugar a la velocidad a la que se producen.
Se empleará una exploración visual para realizar su tratamiento. 

Cabe destacar que a todos los conjuntos de datos se les ha aplicado un filtro de 
tendencias (en este caso el filtro de Hodrick-Prescott \cite{misc:wikipediaHodrick2023}), 
puesto que no existe un interés específico en los valores concretos, sino en las posibles 
tendencias de las diferentes variables, que nos permitan obtener una predicción adecuada 
en cada circunstancia.
De esta manera, se calcula la tendencia diaria de cada periodo temporal.

\section{Valores faltantes}
En cuanto a los valores faltantes, la decisión fue eliminar las muestras que tuvieran
nulos en alguna de sus columnas, puesto que la recuperación empleando técnicas
como la imputación de valores era prácticamente inviable y, por otro lado, se contaban
con ejemplos suficientes como para poder suprimir parte de estos.

\section{Valores extremos}
En cuanto a la detección y el tratamiento de los valores extremos, se empleó el 
rango intercuartílico para realizar la primera de las tareas, sustituyendo las lecturas
del atributo concreto de los ejemplos detectados empleando la mediana diaria.

\imagen{iqr.png}{Representación gráfica del método del rango intercuartílico}{0.75}

De esta manera se eliminaba el ruido presente en alguno de los sensores, a la par 
que no se producían cambios bruscos en el conjunto de datos gracias a emplear la 
mediana del grupo.

\section{Valores erróneos}
En lo referente a los valores erróneos se trata de una de las tareas más arduas 
del proyecto, puesto que parte de esta detección se desarrolló de forma visual directamente
sobre el conjunto de datos de los diferentes sensores y el pluviómetro, siendo, 
de esta manera, en su mayoría un proceso manual.

Para tal propósito se estableció una columna adicional al conjunto de datos original
que indicaba la validez de la muestra concreta.

\section{Visualización de los datos}
En esta sección se explicará el tratamiento realizado en cada uno de los sensores, y,
de esta forma, poner en contexto cómo se han resuelto los problemas detectados
en la fase anterior.

\subsection{Sensor 1}
% Sensor 1
\imagen{processed_data/sensor1.png}{Datos procesados: sensor 1}{1}

Como puede observarse en la Figura \ref{fig:processed_data/sensor1.png}, 
se ha tratado el ruido presente en el conjunto original, sustituyendo los 
respectivos valores por la mediana diaria.

Por otro lado, se ha determinado la supresión de las muestras que indicaban una 
variabilidad brusca de las humedades del suelo debidas a el agotamiento de la batería.
Esta decisión estuvo motivada por la pérdida moderada de datos, además de la 
falta de los valores siguientes.

\newpage

\subsection{Sensor 2}
% Sensor 2
\imagen{processed_data/sensor2.png}{Datos procesados: sensor 2}{1}

En este caso, como puede observarse en la Figura \ref{fig:processed_data/sensor2.png},
se aplicó de nuevo la detección y tratamiento de ruido empleando la mediana como 
valor sustitutorio, además, ante la aparente imposibilidad de obtener una adecuada
recuperación de los datos con alta variabilidad, se optó por no utilizarlos.

\newpage

\subsection{Sensor 3}
% Sensor 3
\imagen{processed_data/sensor3.png}{Datos procesados: sensor 3}{1}

Como puede observarse en la Figura \ref{fig:processed_data/sensor3.png}, en este caso
concreto, el conjunto de datos se modifica, en cuanto a la invalidez de ejemplos se refiere,
visiblemente menos que los sensores anteriores y los siguientes.

En este caso solo fue necesario retirar una porción moderada de las
muestras totales que habían resultado del agotamiento de la batería del sensor.

\newpage

\subsection{Sensor 4}
% Sensor 4
\imagen{processed_data/sensor4.png}{Datos procesados: sensor 4}{1}

Como refleja la Figura \ref{fig:processed_data/sensor4.png}, se optó
por la eliminación de los datos a partir del comienzo de las lecturas que ya
se conocían que eran incorrectas, lo que dejaba un conjunto preparado
para ser empleado en los modelos. 

\newpage

\subsection{Sensor 5}
% Sensor 5
\imagen{processed_data/sensor5.png}{Datos procesados: sensor 5}{1}

En este caso, a diferencia del resto de sensores, no fue necesario una modificación
tan exhaustiva del conjunto de datos, puesto que los únicos tratamientos realizados fueron la 
detección de \textit{outliers} o valores extremos y la eliminación de los ejemplos
con valores faltantes.

\newpage

\subsection{Sensor 6}
% Sensor 6
\imagen{processed_data/sensor6.png}{Datos procesados: sensor 6}{1}

En el caso del sensor 6, como puede apreciarse en la Figura 
\ref{fig:processed_data/sensor6.png}, se ha tratado y eliminado el ruido presente 
mediante la detección de valores extremos. Además, la anomalía de la humedad ha sido
eliminada al tratarse de una situación que no podía corregirse, por ejemplo, con los
métodos mencionados.

\newpage

\subsection{Sensor 7}
% Sensor 7
\imagen{processed_data/sensor7.png}{Datos procesados: sensor 7}{1}

En este caso existían dos problemas principales en el conjunto de datos original:
las sondas de temperatura y humedad comenzaban a realizar lecturas con variaciones
muy elevadas para su profundidad en cortos periodos de tiempo y a la par
había atributos que no contaban con valores, además, al igual que en otros casos,
el agotamiento de la batería producía sondeos incorrectos en las humedades.

Ante estos problemas la solución tomada fue eliminar los ejemplos 
afectados. En el primer caso puesto se trataban de periodos de tiempo moderados y 
en el segundo porque se había producido el desplazamiento del sensor de su posición 
original.

\newpage

\subsection{Sensor 8}
% Sensor 8
\imagen{processed_data/sensor8.png}{Datos procesados: sensor 8}{1}

Como puede observarse en la Figura \ref{fig:processed_data/sensor8.png}, en este caso
no ha sido posible recuperar los datos de gran parte del periodo muestreado.
Los aumentos de humedad mostrados están contrastados como 
provocados por las altas precipitaciones.

Hablando en términos de sesgos, es posible que, al realizar la recuperación de los
valores en los sensores mencionados, puedan haberse sesgado en cierta medida
los conjuntos de datos, sin embargo, estos no dejarán de estar más cercanos a los
valores reales y, por tanto, contener un error más reducido que las muestras originales
con esta casuística.
\apendice{Modelado}

\section{Introducción}
En la fase del modelado, como su nombre indica, se crearán los modelos, de forma que al ser
este proceso (\textit{KDD}) iterativo, se podrá considerar con la etapa anterior, pues existen
herramientas que aunan estos procedimientos.

En este caso, se toman por separado, teniendo en cuenta que los resultados en una
fase afecta de forma directa al resto.

Se estableció el desarrollo de tres redes neuronales diferentes, \textit{MLP}, 
\textit{LSTM} y \textit{GRU}, tratándose las dos últimas de modelos recurrentes.

En esta sección se explicará la transformación del conjunto de datos para poder
realizar el entrenamiento y las respectivas predicciones, así como
las diferentes características de cada uno de los modelos.

\section{Transformación del conjunto de datos}
El primer problema que es necesario afrontar es que estos deberán aceptar varias
entradas de ejemplares, es decir, se necesitarán \(n\) muestras sucesivas para realizar la
regresión.

Por otro lado, también deben ser capaces de realizar \(m\) predicciones, de forma que 
es preciso crear modelos con \(n\) muestras de entrada y \(m\) salidas diferentes.

El segundo problema deriva de la falta de datos en instantes concretos por el procesamiento
realizado previamente, tanto por la supresión de muestras con valores erróneos, como
la de ejemplares con valores faltantes.

Por tanto, se necesita recorrer los datos en busca de los saltos temporales, para obtener
los subconjuntos completos, y de esta manera, asegurar que dentro de éstos estas situaciones
no se producen.

\imagen{window.png}{Selección de la serie de tiempo}{0.75}

En la Fiugura \ref{fig:window.png} puede observarse el proceso esquematizado.
De esta forma, cada uno de los subconjuntos que no tienen saltos, se recorrerá 
mediente ventanas deslizantes formando los datos de entrada (\textit{inputs}) a los 
modelos y las salidas (\textit{targets}).

Cabe destacar que tanto los subconjuntos como las ventanas deberán contener muestras 
suficientes, es decir, \(n + m\) ejemplos.

\newpage

\section{GRU}
En el caso de \textit{GRU}, las celdas de las que se encuentra compuesto el modelo 
emplearán como funciones de activación la tangente hiperbólica y funciones sigmoides, 
es decir, se emplearán los modelos por defecto como puede observars en la Figura 
\ref{fig:gru_cell.png}.

\imagen{gru_cell.png}{Representación celda GRU. Extraido de \cite{wiki:gru2022}}{0.70}

\section{LSTM}
Al igual que en \textit{GRU}, \textit{LSTM} empleará funciones sigmoides y la 
tangente hiperbólica como funciones de activación en las compuertas de las celdas (Figura \ref{fig:LSTM_cell.png}),
empleando, de igual manera, modelos por defecto.

\imagen{LSTM_cell.png}{Representación celda LSTM. Extraido de \cite{wiki:lstm2022}}{0.70}

A diferencia de \textit{GRU}, \textit{LSTM} cuenta con 3 puertas diferentes en cada celda, 
la compuerta de olvido (\textit{Forget Gate}), de entrada (\textit{Input Gate}), 
de salida (\textit{Output Gate}) y un estado oculto (\textit{Cell Satate}), mientras que 
el primer modelo únicamente dispone de dos (\textit{Reset Gate} y \textit{Update Gate}).

De esta forma, en los modelos \textit{LSTM}, se necesitarán realizar más operaciones
que en \textit{GRU}, además de almacenar más información por celda, 
de manera que el tiempo de entrenamiento será generalmente mayor.

\section{MLP}
El perceptrón multicapa, a diferencia de los otros dos modelos, no acepta
la entrada de secuencias variables, por lo que se debe establecer una neurona de entrada
por cada uno de los atributos, es decir, la red contendrá
\(a * n\) neuronas de entrada, siendo \(a\) el número de características de los ejemplos y
\(n\) el número de muestras de entrada.

De forma similar, contará con \(a * m\) número de neuronas en la capa de salida, siendo
\(m\) el número de ejemplares de salida.

\imagen{mlp_anexos.png}{Diagrama de entradas y salidas MLP}{0.75}

Se establece una función \textit{ReLU} (Figura \ref{fig:relu.png}) como función
de activación en las neuronas del modelo.

\imagen{relu.png}{Función ReLU}{0.70}

A pesar de que \textit{MLP} no acepta entradas secuenciales y que por tanto requiere de
tantas neuronas de entrada como número de valores hay en la secuencia de entrada,
el tiempo de entrenamiento de estos modelos será menor generalmente que el de los
recurrentes, puesto que se realizan menos operaciones por término medio,
además de necesitar de una cantidad de memoria menor para almacenar los pesos sinápticos.

\apendice{Evaluación}

\section{Introducción}
En esta fase se estima el rendimiento de los diferentes modelos, y, en su caso, 
se reconsideran los objetivos del proyecto, de manera que, si estos son poco efectivos,
se vuelve a la primera fase.

En todas las redes neuronales se ha empleado validación cruzada externa y, además, se realizarán
diferentes ejecuciones con diversos parámetros para conocer el rendimiento de estas, a 
la par que poder eliminar en cierta medida el efecto del azar en los modelos 
resultantes.

Para validarlos se realizarán pruebas variando el número de capas
ocultas, así como el de celdas/neuronas en cada una de estas y el ratio de aprendizaje.
De esta manera, se realizarán combinaciones con 1, 2 y 3 capas en cada uno de los modelos, 
así como 32, 16 y 8 celdas/neuronas y 0.1, 0.01 y 0.001 de ratio de aprendizaje. 

Por otro lado, se ha dividido el conjunto de datos original en un conjunto de entrenamiento del 75\% de los disponibles
por sensor, 15\% para validación siendo el restante de test.

Con respecto al número de muestras de entrada se emplearán 6 horas previas para realizar
una predicción, empleando un tamaño de bloque de 256 ejemplos (los modelos 
no se entrenan con todos los datos en todas las iteraciones, sino que se emplean un subconjunto).
En la salida, por otro lado, se realizarán predicciones de 1, 2 y 3 horas (marcadas como MSE 1, MSE 2 y MSE 3 en las
respectivas tablas).

Además, se emplea un planificador de la tasa de aprendizaje, de forma
que en la época 50 disminuye multiplicando \(lr * e^{-0.1}\), para, de esta manera,
centrar los esfuerzos en la región explorada.

Para cada una de las pruebas se realizarán 5 ejecuciones diferentes con 100 iteraciones, 
tomando como medida comparativa el error cuadrático medio de validación, lo que supone 405 ejecuciones por modelo, es decir, 1215 ejecuciones en total.

\newpage

\section{Resultados de GRU}
En la Tabla \ref{tabla:Evaluación GRU} se presentan los errores cuadráticos medios
que se obtendrán para cada una de las pruebas. En esta se puede observar los resultados obtenidos con 1, 2 y 3 predicciones 
como salida de los modelos. 

\tablaSmallSinColores{Evaluación GRU}{l l l l l l}{Evaluación GRU}{\textbf{Capas} & \textbf{Tamaño capa} & \textbf{LR} & \textbf{MSE 1} & \textbf{MSE 2} & \textbf{MSE 3}\\}
{
    1 & 8 & 0.001 & 1.72183E-05 & 3.98320E-05 & 5.84809E-05 \\
    1 & 8 & 0.01 & 3.37952E-06 & 7.99268E-06 & 1.41913E-05\\
    1 & 8 & 0.1 & 0.00551 & 0.00127 & 0.00188 \\
    1 & 16 & 0.001 & 9.1907E-06 & 1.73932E-05 & 3.24371E-05 \\
    1 & 16 & 0.01 & 9.55771E-07 & 2.96765E-06 & 7.32603E-06 \\
    1 & 16 & 0.1 & 0.00474 & 0.00718 & 0.00138 \\
    1 & 32 & 0.001 & 4.50204E-06 & 1.22795E-05 & 2.51439E-05 \\
    1 & 32 & 0.01 & \textbf{5.96471E-07} & 1.69264E-06 & 4.50310E-06 \\
    1 & 32 & 0.1 & 0.00297 & 0.00158 & 0.00175 \\

    % 2 capas ocultas
    2 & 8 & 0.001 & 1.56769E-05 & 2.90745E-05 & 4.93188E-05 \\
    2 & 8 & 0.01 & 1.65277E-06 & 6.30449E-06 & 1.14688E-05 \\
    2 & 8 & 0.1 & 0.00965 & 0.00972 & 0.00026 \\
    2 & 16 & 0.001 & 5.07987E-06 & 1.33184E-05 & 2.59817E-05 \\
    2 & 16 & 0.01 & \textbf{5.66822E-07} & \textbf{1.45485E-06} & \textbf{3.49998E-06} \\
    2 & 16 & 0.1 & 0.01148 & 0.01579 & 0.00363 \\
    2 & 32 & 0.001 & 2.18717E-06 & 7.23831E-06 & 1.42593E-05 \\
    2 & 32 & 0.01 & \textbf{2.69205E-07} & \textbf{9.00049E-07} & \textbf{4.04636E-06} \\
    2 & 32 & 0.1 & 0.01727 & 0.00611 & 0.00797 \\

    % 3 capas ocultas
    3 & 8 & 0.001 & 1.61179E-05 & 2.57331E-05 & 4.20175E-05 \\
    3 & 8 & 0.01 & 1.42658E-06 & 5.39812E-06 & 9.51720E-06 \\
    3 & 8 & 0.1 & 0.00706 & 0.00029 & 0.00242 \\
    3 & 16 & 0.001 & 5.18860E-06 & 1.23398E-05 & 2.02378E-05 \\
    3 & 16 & 0.01 & \textbf{4.66034E-07} & \textbf{1.29760E-06} & \textbf{3.49093E-06}\\
    3 & 16 & 0.1 & 0.02418 & 0.01457 & 0.01497\\
    3 & 32 & 0.001 & 2.36664E-06 & 5.75822E-06 & 1.13415E-05\\
    3 & 32 & 0.01 & 8.90246E-07 & \textbf{9.57532E-07} & \textbf{2.41350E-06}\\
    3 & 32 & 0.1 & 0.00976 & 0.01064 & 0.01097\\
}

Como puede observarse, las mayores eficiencias en términos generales se obtienen 
con un \textit{learning rate} de 0.01, destacando los modelos de 2 capas con 
32 y 16 celdas de memoria respectivamente.

Cabe destacar, además, que, como se podía intuir inicialmente, los errores aumentan generalmente
a medida que realizamos más predicciones como salida.

A continuación se presentan las comparativas de los errores cuadráticos medios promedio para un número
diferente de predicciones (1, 2 y 3 predicciones) del modelo estudiado y un parámetro diferente en cada
ocasión.

En el caso del número de capas, como puede observarse en la Figura \ref{fig:Capas_GRU.png}, en promedio
se obtendrán modelos con menor error generalmente en sistemas con 3 predicciones a la salida, siendo los 
modelos de 1 capa los que mayor precisión han presentado en las pruebas.
\imagen{Capas_GRU.png}{Comparativa de número de capas para GRU}{0.9}

En lo que se refiere a las celdas de memoria, por término medio generalmente se obtendrán modelos más
precisos con estructuras que realizan 3 predicciones por entrada, siendo la configuración de 8 unidades 
la que proporciona modelos con errores menores en promedio.
\imagen{Unidades_GRU.png}{Comparativa de número de unidades para GRU}{0.9}

Como puede apreciarse en la Figura \ref{fig:Lr_GRU.png}, con 0.1 como ratio de aprendizaje se obtienen los
resultados con mayor error en promedio.
\imagen{Lr_GRU.png}{Comparativa de ratio de aprendizaje para GRU}{0.9}

En la Figura \ref{fig:Lr_GRU_ampl.png} puede observarse la diferencia del resto de los valores del 
parámetro. En este caso, se obtienen mejores resultados por término medio con un ratio de aprendizaje de
0.01, siendo los modelos de 1 predicción por entrada los que generalmente obtienen errores menores.
\imagen{Lr_GRU_ampl.png}{Comparativa de ratio de aprendizaje para GRU (0.01 y 0.001)}{0.9}

\newpage

\section{Resultados de LSTM}
Al igual que en la situación anterior, se realizarán pruebas con el mismo número de predicciones
por cada entrada, obteniéndose los resultados reflejados en la Tabla \ref{tabla:Evaluación LSTM}.

\tablaSmallSinColores{Evaluación LSTM}{l l l l l l}{Evaluación LSTM}{\textbf{Capas} & \textbf{Tamaño capa} & \textbf{LR} & \textbf{MSE 1} & \textbf{MSE 2} & \textbf{MSE 3}\\}
{
    1 & 8 & 0.001 & 2.12520E-05 & 3.03812E-05 & 5.60944E-05 \\
    1 & 8 & 0.01 & 3.4556E-06 & 1.05315E-05 & 1.68195E-05 \\
    1 & 8 & 0.1 & 3.9049E-06 & 2.55932E-05 & 2.56167E-05 \\
    1 & 16 & 0.001 & 1.18867E-05 & 1.98377E-05 & 4.04516E-05 \\
    1 & 16 & 0.01 & 1.31837E-06 & 3.46421E-06 & 7.71354E-06 \\
    1 & 16 & 0.1 & 3.14115E-06 & 7.24853E-06 & 7.55877E-05 \\
    1 & 32 & 0.001 & 5.84463E-06 & 1.46677E-05 & 2.83427E-05 \\
    1 & 32 & 0.01 & \textbf{6.62421E-07} & \textbf{1.96414E-06} & 5.22169E-06 \\
    1 & 32 & 0.1 & \textbf{9.51891E-07} & 0.00011 & 1.05361E-05 \\

    % 2 capas ocultas
    2 & 8 & 0.001 & 3.29916E-05 & 4.19003E-05 & 6.14671E-05 \\
    2 & 8 & 0.01 & 4.94574E-06 & 8.51327E-06 & 1.33202E-05 \\
    2 & 8 & 0.1 & 6.85972E-06 & 8.04903E-06 & 1.58783E-05 \\
    2 & 16 & 0.001 & 1.33461E-05 & 2.75447E-05 & 4.1725E-05 \\
    2 & 16 & 0.01 & 1.23340E-06 & \textbf{2.38821E-06} & \textbf{4.63577E-06} \\
    2 & 16 & 0.1 & 6.65626E-06 & 2.42052E-05 & 1.43046E-05 \\
    2 & 32 & 0.001 & 5.63944E-06 & 1.15909E-05 & 2.34114E-05 \\
    2 & 32 & 0.01 & \textbf{4.56945E-07} & \textbf{1.43223E-06} & \textbf{2.78594E-06}\\
    2 & 32 & 0.1 & 4.20038E-06 & 1.00814E-05 & 7.57666E-05\\

    % 3 capas ocultas
    3 & 8 & 0.001 & 4.58400E-06 & 5.81849E-05 & 0.00012 \\
    3 & 8 & 0.01 & 5.90159E-06 & 9.07351E-06 & 1.48364E-05 \\
    3 & 8 & 0.1 & 0.00031 & 0.00486 & 1.73157E-05 \\
    3 & 16 & 0.001 & 2.22344E-05 & 3.47894E-05 & 5.07083E-05 \\
    3 & 16 & 0.01 & 1.17654E-06 & 2.41703E-06 & \textbf{4.06992E-06} \\
    3 & 16 & 0.1 & 0.00487 & 1.02880E-05 & 0.00490 \\
    3 & 32 & 0.001 & 8.16043E-06 & 1.64882E-05 & 2.57956E-05 \\
    3 & 32 & 0.01 & \textbf{6.54021E-07} & \textbf{1.19201E-06} & \textbf{2.58606E-06} \\
    3 & 32 & 0.1 & 0.00491 & 0.01459 & 0.00986 \\
}

De forma similar que en \textit{GRU}, en \textit{LSMT} se generan mejores resultados con un 
\textit{learning rate} de 0.01, destacando, de igual forma, los modelos con 2 capas, además
de los de 3 capas y 32 celdas de memoria.
Al igual que antes, los errores se incrementan, generalmente, con el aumento del número de predicciones.

A continuación, del mismo modo que en la sección anterior, se presentan las comparativas de los errores 
cuadráticos medios promedio para un número diferente de predicciones (1, 2 y 3 predicciones) del modelo estudiado y un parámetro diferente en cada ocasión.

En la Figura \ref{fig:Capas_LSTM.png} se observa que las configuraciones de 2 capas registran los 
errores más bajos en promedio que el resto de valores del parámetro, siendo los modelos de 1 predicción
futura los que generalmente obtendrán mejores resultados.
\imagen{Capas_LSTM.png}{Comparativa de número de capas para LSTM}{0.9}

En el caso del número de celdas de memoria, se obtendrán modelos más precisos por término medio con 8 
unidades por capa (a excepción de los sistemas con 2 predicciones por entrada que contarán con mejor 
desempeño en redes con 16 unidades). 
\imagen{Unidades_LSTM.png}{Comparativa de número de unidades para LSTM}{0.9}

En lo referente al ratio de aprendizaje, los valores del parámetro para los que se obtienen modelos
con menor precisión por término medio será 0.1, seguido de 0.001.
\imagen{Lr_LSTM.png}{Comparativa de ratio de aprendizaje para LSTM}{0.9}

En la Figura \ref{fig:Lr_LSTM_ampl.png} puede observarse que para 0.01 de ratio de aprendizaje se obtienen
los resultados con mayor precisión por término con 1 y 2 predicciones por entrada. Siendo, de esta forma,
la configuración de 1 predicción la que menores errores arroje en promedio en comparación con el resto
de parámetros.
\imagen{Lr_LSTM_ampl.png}{Comparativa de ratio de aprendizaje para LSTM (0.01)}{0.9}

\newpage

\section{Resultados de MLP}
En la Tabla \ref{tabla:Evaluación MLP} se proporcionan los diferentes resultados obtenidos con \textit{MLP}.

\tablaSmallSinColores{Evaluación MLP}{l l l l l l}{Evaluación MLP}{\textbf{Capas} & \textbf{Tamaño capa} & \textbf{LR} & \textbf{MSE 1} & \textbf{MSE 2} & \textbf{MSE 3} \\}
{
    1 & 8 & 0.001 & 0.00076 & 0.00017 & 0.00089 \\
    1 & 8 & 0.01 & 0.01218 & 0.00313 & 0.00229 \\
    1 & 8 & 0.1 & 0.02064 & 0.02415 & 0.02420 \\
    1 & 16 & 0.001 & 6.21736E-05 & 7.85435E-05 & 6.44446E-05 \\
    1 & 16 & 0.01 & 0.00059 & 0.00119 & 8.35920E-05 \\
    1 & 16 & 0.1 & 0.02056 & 0.02415 & 0.02420 \\
    1 & 32 & 0.001 & 1.33441E-05 & \textbf{2.31054E-05} & 4.27180E-05 \\
    1 & 32 & 0.01 & \textbf{5.57840E-06} & \textbf{1.32064E-05} & \textbf{2.49582E-05} \\
    1 & 32 & 0.1 & 0.02415E-05 & 0.02062 & 0.02420 \\

    % 2 capas ocultas
    2 & 8 & 0.001 & 0.00094 & 0.00043 & 0.00085 \\
    2 & 8 & 0.01 & 0.00673 & 0.00332 & 0.00327 \\
    2 & 8 & 0.1 & 0.02428 & 0.02430 & 0.02063 \\
    2 & 16 & 0.001 & 1.96797E-05 & 2.97114E-05 & 0.00013 \\
    2 & 16 & 0.01 & 0.00042 & 0.00040 & 0.00089 \\
    2 & 16 & 0.1 & 0.01624 & 0.02422 & 0.01343 \\
    2 & 32 & 0.001 & 8.32571E-06 & \textbf{2.15010E-05} & 3.63325E-05 \\
    2 & 32 & 0.01 & \textbf{5.56436E-06} & 5.32363E-05 & \textbf{2.65348E-05} \\
    2 & 32 & 0.1 & 0.02032 & 0.02060 & 0.02043 \\

    % 3 capas ocultas
    3 & 8 & 0.001 & 0.00152 & 0.00283 & 0.00195 \\
    3 & 8 & 0.01 & 0.00168 & 0.00284 & 0.00188 \\
    3 & 8 & 0.1 & 0.02417 & 0.02428 & 0.02041 \\
    3 & 16 & 0.001 & 5.74476E-05 & 5.53494E-05 & 0.00016 \\
    3 & 16 & 0.01 & 0.00031 & 0.00015 & 0.00095 \\
    3 & 16 & 0.1 & 0.02035 & 0.02036 & 0.01269 \\
    3 & 32 & 0.001 & \textbf{6.49630E-06} & \textbf{2.14812E-05} & \textbf{3.33377E-05} \\
    3 & 32 & 0.01 & \textbf{6.91262E-06} & 3.07683E-05 & \textbf{2.54072E-05} \\
    3 & 32 & 0.1 & 0.02028 & 0.01674 & 0.01653 \\
}

Al igual que en las anteriores secciones, se presentan las comparativas de los errores 
cuadráticos medios promedio para un número diferente de predicciones (1, 2 y 3 predicciones) del modelo 
estudiado y un parámetro diferente en cada ocasión.

En el caso del número de capas, se puede observar (Figura \ref{fig:Capas_MLP.png}) que a diferencia de
los otros modelos, la influencia del número de capas en la precisión es aparentemente menor.
\imagen{Capas_MLP.png}{Comparativa de número de capas para MLP}{0.9}

En cuanto al número de neuronas por capa, puede apreciarse que con 32 unidades se obtienen los modelos
con mayor precisión por término medio, siendo la configuración de 2 predicciones por entrada la que
generalmente arrojará resultados con menor error.
\imagen{Unidades_MLP.png}{Comparativa de número de unidades para MLP}{0.9}

Como puede observarse en la Figura \ref{fig:Lr_MLP.png}, el ratio de aprendizaje tiene una aparente
mayor influencia sobre los resultados, siendo la configuración de 0.001 la que menores errores por
término medio proporcionan.
\imagen{Lr_MLP.png}{Comparativa de ratio de aprendizaje para MLP}{0.9}

\section{Comparativa entre modelos}
\imagen{Capas_Comp.png}{Comparativa de número de capas}{0.9}
\imagen{Unidades_Comp.png}{Comparativa de número de unidades}{0.9}
\imagen{Lr_Comp.png}{Comparativa de ratio de aprendizaje}{0.9}
\apendice{Documentación técnica de programación}

\section{Introducción}
En esta sección se incluyen la documentación técnica del programador, presentando la estructura de directorios del proyecto, junto con el manual para realizar la correcta 
instalación y ejecución de este. 

\section{Estructura de directorios}
El proyecto cuenta con la siguiente estructura de directorios:
\begin{itemize}
    \item \textbf{/data/:} directorio que contiene los diferentes datos del proyecto, tanto procesados, sin procesar y los datos integrados que se emplearán en el modelado.
    \subitem \textbf{/data/raw/:} directorio con los datos sin procesar (úncamente con la selección previa de validez).
    \subitem \textbf{/data/processed/:} directorio con los datos procesados.
    \subitem \textbf{/data/integrated/:} directorio con los datos integrados en un único fichero.
    \item \textbf{/img/graphics/:} directorio con las diferentes gráficas resultado de la ejecución de los scripts de graficado.
    \item \textbf{/scripts/:} directorio con los scripts para la instalación de los entornos virtuales de Python junto con los requerimientos para ejecutar todos los ficheros
        fuente del proyecto.
    \item \textbf{/src/:} directorio con los diferentes ficheros fuente y variables de entorno y globales.
    \item \textbf{/models/:} directorio con los diferentes modelos obtenidos en el proceso final. Uno subdirectorio para cada diferente modelo neuronal implementado.
\end{itemize}

\section{Manual del programador}
En esta subsección se explicará cómo realizar una correcta descarga e instalación de los entornos necesarios para llevar a cabo la ejecución del proyecto.

Para descargar todo el contenido es necesario tener instalado en el sistema \textbf{Git}. Es posible clonar el repositorio introduciendo en la consola de 
git: 
\begin{lstlisting}[language=Bash]
    git clone https://github.com/GabiHV/TFG22-23
\end{lstlisting}

De igual forma, para poder llevar a cabo la ejecución e instalación del resto de las dependencias es necesario tener instalado Python 3.9.13.

Para instalar el intérprete del lenguaje empleado en el proyecto se debe acudir a la página web oficial de los desarrolladores y descargar el ejecutable de instalación oficial.
La instalación puede realizarse en el siguiente enlace \cite{misc:python2023}.
En la fuente mencionada se pueden escoger diferentes formas de instalación. Dependiendo del sistema operativo utilizado en la máquina en la que se ejecutará el proyecto se debe seleccionar una
u otra y seguir los pasos establecidos.

Durante el desarrollo del proyecto se empleó como entorno de programación Visual Studio Code \cite{misc:code2023}, sin embargo para su ejecución podemos emplear otros entornos como 
Anaconda Navigator \cite{misc:conda2023}.
Se explicará los pasos a seguir con el editor mencionado para habilitar los entornos, puesto que simplifica el trabajo al disponer de scripts que realizan de forma automática 
la instalación de las dependencias.
Los ficheros mencionados se encuentran en el directorio \textbf{/scripts/}.

Para ejecutar el script correspondiente al entorno de PowerShell de Windows se necesita establecer la política que lo permita. Para ello se debe abrir la terminal mencionada
como administrador del sistema e introducir: 
\begin{lstlisting}[language=Bash]
    Set-ExecutionPolicy Unrestricted    
\end{lstlisting}

Tras esto, se puede introducir el siguiente comando para iniciar el proceso de instalación: 
\begin{lstlisting}
    ./Virtual_env.ps1
\end{lstlisting}

Para ejecutar el script en el CMD de Windows se introduce:
\begin{lstlisting}[language=Bash]
    virtual_env.bat
\end{lstlisting}

De forma similar en Linux Bash:
\begin{lstlisting}[language=Bash]
    chmod +x virtual_env.sh && ./virtual_env.sh    
\end{lstlisting}

Una vez finalice el proceso de instalación de todas las dependencias se podrán ejecutar los diferentes ficheros fuente de Python Notebook abriendo el proyecto en Visual Studio Code
y estableciendo el Kernel de ejecución al entorno configurado. Está definido que el entorno virtual se denomine \textbf{.venv}, por lo que será necesario buscar entre los diferentes 
instalados haciendo click en la parte superior derecha del notebook (en el botón para la selección del intérprete de ejecución).

\imagen{venv_selection.png}{Búsqueda del entorno virtual}{1}
\imagen{venv_selected.png}{Selección del entorno virtual}{0.25}

Posteriormente podrá ejecutarse cualquier fichero Pyhton Notebook con el botón de ``Execute All'', o proceder por celda.

En lo referente a los ficheros fuente de Python, con el entorno virtual instalado abriendo una consola en el sistema operativo en \textbf{/.venv/bin/} en Linux y \textbf{/.venv/Scripts/}
en Windows, podremos ejecutar la activación del entorno virutal con los scripts incluidos en el directorio (ejecutando \linebreak \textbf{activate.bat} y \textbf{Activate.ps1} en Windows dependiendo del terminal empleado).
Esta operación se realizará automáticamente al realizar la instalación de los módulos de Python incluidos en el fichero de requerimientos, por lo que se puede aprovechar la terminal en ejecución para 
este propósito.

Para ejecutar los ficheros de Python adecuadamente es necesario estar situados con la terminal en el directorio correspondiente. En este caso se debe nevegar hasta \textbf{./src/}

\section{Pruebas del sistema}
En esta subsección se presentará la forma de realizar las modificaciones en los hiperparámetros de los modelos, de forma que estos puedan variar de acuerdo a los nuevos requerimientos introducidos.

Los modelos contienen los siguientes hiperparámetros:
\begin{itemize}
    \item \textbf{learning\_rate:} ratio de aprendizaje empleado en la variación de los pesos en los modelos neuronales.
    \item \textbf{batch\_size:} tamaño del conjunto de datos que se emplea en una única iteración en el proceso
        de aprendizaje.
    \item \textbf{epochs:} cantidad de épocas que se entrenará cada modelo.
    \item \textbf{window\_size\_inputs:} tamaño de la ventana de datos que se introduce como datos de entrada al modelo
        (se corresponde con el número de horas previas para realizar X predicciones).
    \item \textbf{window\_size\_targets:} tamaño de la ventana de datos que se emplean como datos a predecir.
    \item \textbf{train\_frac:} fracción del conjunto total de datos que se empleará para entrenar los modelos.
    \item \textbf{val\_frac:} fracción del conjunto total de datos que se empleará para validar los modelos, siendo el 
        \(1 - train\_frac - val\_frac\) la fracción del conjunto de test.
\end{itemize}

Por cada uno de los diferentes modelos se proporcionará una gráfica del error de entrenamiento y validación
durante el proceso correspondiente, además de las gráficas comparativas de los valores predichos y 
los reales por sensor y atributo, así como el error máximo total en cada característica para todos los sensores y las gráficas
de los errores en el conjunto de test, de forma que se 
pueda dar una idea de los valores en los que ronda el error en cada una de las variables. 


\apendice{Documentación de usuario}

\section{Introducción}
En esta sección se presentará la manera de cargar los modelos resultantes en un fichero Python para poder ser desplegados en un producto software, 
así como la forma que deberá tener el tensor de entrada a la red y los datos de salida.

\section{Requisitos de usuarios}
En cuanto a los requisitos del usuario, será necesario tener instalado \textbf{Python 3.9.13} \cite{misc:python2023}, junto con la versión 2.11.0 de \textbf{TensorFlow} \cite{misc:tensorflow2023},
de manera que se puedan cargar los modelos almacenados empleando la función de la \textit{API} \textbf{Keras} para tal propósito.

\section{Instalación}
En cuanto a la instalación de la versión concreta del intérprete de Python puede realizarse en \cite{misc:python2023}, mientras que, para instalar la dependencia de la \textit{API},
se puede introducir el comando:
\begin{lstlisting}[language=Bash]
    pip3 install tensorflow=2.11.0
\end{lstlisting}

Para cargar un modelo con la librería mencionada, se debe emplear \cite{misc:tensorflow_save2023}:
\begin{lstlisting}[language=Python]
    from tensorflow import keras
    keras.models.load_model('<path_del_modelo>')
\end{lstlisting}

\section{Manual del usuario}
Para realizar predicciones con el modelo pertinente, las entradas deben tener una estructura concreta que dependerá de cómo se haya entrenado a cada una de las diferentes redes neuronales.
Es decir, el tamaño del número de muestras de entrada dependerá de la cantidad de ``\textit{backtracking}'' que se haya establecido en el entrenamiento.

En este caso se debe introducir un tensor bidimensional de 6 muestras (se han obtenido modelos que aceptan 6 horas) con 7 variables de entrada cada una que se corresponde a:
\begin{itemize}
    \item \textbf{t\_ext:} temperatura exterior media en una hora (-50, 50).
    \item \textbf{h\_ext:} humedad exterior media en una hora (0, 100).
    \item \textbf{t\_C\_cal:} temperatura media de la sonda de temperatura más superficial en una hora (-50, 50).
    \item \textbf{h\_C\_cal:} humedad media de la sonda de humedad más superficial en una hora (0, 100).
    \item \textbf{t\_L\_cal:} temperatura media de la sonda de temperatura interna en una hora (-50, 50).
    \item \textbf{h\_L\_cal:} humedad media de la sonda de humedad interna en una hora (0, 100).
    \item \textbf{sensor:} sensor al que se corresponde los datos (0, 7).
\end{itemize}

Por otro lado, las variables deberán estar normalizadas en el rango 0-1, con los máximos y mínimos especificados anteriormente.

En el caso del modelo \textit{MLP}, en lugar de un tensor tridimensional, se deberá modificar para que se corresponda con un vector del número de variables por el de muestras.

La salida será igualmente un tensor bidimensional del número de predicciones establecidas en el entrenamiento del modelo con 6 variables cada una que se corresponden a las mencionadas
anteriormente a excepción del número de sensor. 

De forma inversa, el modelo proporcionará datos normalizados, por lo que para obtener cada atributo en un rango correcto deberá de denormalizarse.


\bibliographystyle{IEEEtran}
\bibliography{bibliografiaAnexos}

\end{document}
